\documentclass[11pt]{article}
\usepackage{amsmath,amssymb,amsthm,mathtools}
\usepackage[margin=1in]{geometry}
\theoremstyle{plain}
\newtheorem{theorem}{Theorem}
\newtheorem{lemma}{Lemma}
\theoremstyle{remark}
\newtheorem{remark}{Remark}
\newcommand{\E}{\mathbb{E}}\newcommand{\R}{\mathbb{R}}
\newcommand{\dd}{\,\mathrm d}

\title{Formalization brief: Karamata's Tauberian theorem at constant $L$\\
(the case the project uses), with a fidelity repair of the general statement}
\date{7 July 2026}

\begin{document}
\maketitle

\noindent\textbf{Goal.} In \texttt{TypeDDecouplingLCLT.lean} (attached),
replace the \texttt{sorry}-ed \texttt{karamata\_tauberian} by a PROVED theorem
at the strength the project uses, and rewire its sole consumer
\texttt{lem\_tau}. Abstract analysis in a new library-clean file
\texttt{TypeDDecouplingKaramata.lean} encouraged. Only these may change; whole
project must build; standard axioms only.

\textbf{Fidelity repair (document as such).} The encoded general statement
quantifies over ALL pointwise slowly varying $L$ with no measurability
hypothesis; the cited theorem (BGT Thm 1.7.1$'$, Feller XIII.5) requires $L$
measurable, and the slowly-varying theory (UCT) fails for non-measurable $L$
--- the encoded statement therefore exceeds the cited literature and is not a
faithful citation. Repair: restate the theorem at constant $L$ (which is what
the consumer \texttt{lem\_tau} instantiates: $L\equiv m(r)/(2\sqrt a)$,
$\rho=1/2$), i.e.
\[
  \omega(\lambda)=\int_0^\infty e^{-\lambda t}p(t)\dd t,\qquad
  \omega(\lambda)\sim c\,\lambda^{-\rho}\ (\lambda\downarrow0)
  \iff
  \int_0^sp\sim \frac{c\,s^{\rho}}{\Gamma(\rho+1)}\ (s\to\infty),
\]
for $p\ge0$ locally integrable, $c>0$, $\rho\ge0$. Keep a docstring recording:
the general measurable-$L$ theorem is the cited literature; the constant-$L$
case is proved here; the removed statement's measurability gap (the sixth
fidelity catch of this project). Both directions ($\iff$) are wanted; if one
direction resists, the Tauberian direction ($\Rightarrow$ from $\omega$ to
$\int p$) is the one \texttt{lem\_tau} consumes and is the minimum
deliverable (then keep the Abelian direction as a separate statement or drop
the iff, documenting the choice).

\section*{Tier 1: the Abelian direction (elementary)}

With $U(s):=\int_0^sp$: Fubini on $\{(t,s):0<t\le s\}$ gives
$\omega(\lambda)=\lambda\int_0^\infty e^{-\lambda s}U(s)\dd s$; substitute
$s=x/\lambda$:
$\lambda^{\rho}\omega(\lambda)=\int_0^\infty e^{-x}\,\lambda^{\rho}
U(x/\lambda)\dd x$. From $U(s)\sim cs^\rho/\Gamma(\rho+1)$: pointwise
$\lambda^\rho U(x/\lambda)\to cx^\rho/\Gamma(\rho+1)$, dominated by
$C(1+x^\rho)$ (from the equivalence for large $s$ and monotonicity of $U$ for
the head), so DCT gives
$\lambda^\rho\omega(\lambda)\to\frac c{\Gamma(\rho+1)}\int_0^\infty
x^\rho e^{-x}\dd x=c$.

\section*{Tier 2: the Tauberian direction (Karamata's method, constant $L$)}

Assume $\lambda^\rho\omega(\lambda)\to c$. Handle $\rho=0$ separately
(monotone convergence: $\omega(\lambda)\uparrow U(\infty)$ as
$\lambda\downarrow0$ by MCT for $p\ge0$, so $U(\infty)=c$ finite and the
conclusion is $U(s)\to c$). For $\rho>0$:
\begin{itemize}
\item[(a)] \underline{Moments.} For each $k\in\mathbb N$,
$\lambda^\rho\,\omega\big((k{+}1)\lambda\big)\to c$ gives
\[
  \int_0^\infty e^{-(k+1)\lambda t}\lambda^{\rho}p(t)\dd t
  \ \longrightarrow\ \frac{c}{(k+1)^{\rho}}
  \ =\ \frac c{\Gamma(\rho)}\int_0^\infty e^{-(k+1)u}u^{\rho-1}\dd u .
\]
Hence for every polynomial $P$ (in the variable $e^{-u}$),
$\int_0^\infty P(e^{-\lambda t})e^{-\lambda t}\lambda^{\rho}p(t)\dd t\to
\frac c{\Gamma(\rho)}\int_0^\infty P(e^{-u})e^{-u}u^{\rho-1}\dd u$
(finite linear combinations).
\item[(b)] \underline{One-sided approximation.} The target is
$g_0(u):=\mathbf 1_{[0,1]}(u)$, since
$\int_0^\infty\mathbf 1_{[0,1]}(\lambda t)\,\lambda^\rho p(t)\dd t
=\lambda^\rho U(1/\lambda)$ and
$\frac c{\Gamma(\rho)}\int_0^\infty\mathbf 1_{[0,1]}(u)u^{\rho-1}\dd u
=\frac c{\Gamma(\rho+1)}$. Write test functions as $h(e^{-u})$,
$h:[0,1]\to\R$; the target is $h_0(x)=\mathbf 1_{[e^{-1},1]}(x)$. Construct,
for each $\eps>0$, CONTINUOUS $h_-\le h_0\le h_+$ on $[0,1]$ differing from
$h_0$ only on $(e^{-1}-\eps',e^{-1}+\eps')$, then Stone--Weierstrass /
Bernstein (in Mathlib) gives polynomials $P_\pm$ with
$h_-{-}\eps\le P_-\le h_-$ and $h_+\le P_+\le h_+{+}\eps$
(one-sidedness by shifting the uniform approximation by $\eps$). Using
$p\ge0$, sandwich
\[
  \int P_-(e^{-\lambda t})e^{-\lambda t}\cdot\frac{stuff}{}\ \le\
  \lambda^\rho U(1/\lambda)\text{-type}\ \le\ \int P_+(\cdots)
\]
--- CARE: the moment functionals in (a) carry the factor $e^{-u}$ ($k\ge0$,
i.e.\ test functions $x\cdot(\text{poly})$, vanishing at $x{=}0$
$\Leftrightarrow$ $u{=}\infty$); arrange the sandwich for
$h_0(x)/x$-weighted functionals or include the factor in $h_\pm$
(the classical bookkeeping --- either convention works; the limit measure
$\frac c{\Gamma(\rho)}u^{\rho-1}\dd u\circ(e^{-u})^{-1}$ integrates the
approximation error to $O(\eps)$ since the jump point carries no mass and,
for the $x\to0$ end, $u^{\rho-1}e^{-u}$ is integrable).
\item[(c)] \underline{Conclude.} $\limsup/\liminf$ of
$\lambda^\rho U(1/\lambda)$ within $O(\eps)$ of $c/\Gamma(\rho+1)$; let
$\eps\to0$; convert $\lambda\downarrow0$ to $s\to\infty$ ($s=1/\lambda$; the
$\mathrm{IsEquivalent}$ form as in the current statement).
\end{itemize}

\section*{Tier 3: rewiring}

\texttt{lem\_tau} currently instantiates the old statement with a constant
$L$; rewire to the new theorem (the instantiation should simplify). Remove the
old general statement; its docstring content (BGT/Feller citations) moves to
the new theorem's docstring together with the fidelity note. Target count
9 $\to$ 8.

\begin{remark}[hygiene]
(1) Search Mathlib first: Bernstein polynomials
(\texttt{bernsteinApproximation}), Stone--Weierstrass, $\Gamma$-integrals,
Laplace-transform lemmas, \texttt{Asymptotics.IsEquivalent} API.
(2) Sanctioned fallbacks: prove the Tauberian direction only (keep Abelian as
a separate proved lemma or drop, documented); if general $\rho>0$ resists in
(b)'s bookkeeping, $\rho=1/2$ with the explicit weight
$u^{-1/2}e^{-u}$ is the minimum deliverable (what \texttt{lem\_tau} needs).
(3) Constants/forms free; match \texttt{lem\_tau}'s consumption exactly.
(4) Report which directions/generality were achieved and the final count.
\end{remark}

\end{document}
